\section{Introduction}\label{sec:intro}

Automatic Chord Recognition (ACR) underpins computational musicology, structural segmentation, and content-based audio indexing~\cite{pauwels2019twenty}. Although ACR architectures have progressed from handcrafted chroma templates and hidden Markov models to learned deep chroma extractors~\cite{korzeniowski2016feature}, convolutional-recurrent networks~\cite{jiang2019large}, and self-attention Transformers~\cite{park2019btc,gulati2020conformer}, performance on complex chords remains bounded: classifiers recognize major and minor triads far better than four-note tetrads and extended harmonies~\cite{pauwels2019twenty} (e.g., $77.6\%$ MajMin versus $58.8\%$ Tetrads weighted recall for a K-stream HMM on the canonical $1{,}217$-song benchmark~\cite{humphrey2015four}).

This ceiling appears driven less by model capacity than by the severe Zipfian skew of real-world musical corpora: in the $1{,}217$-song benchmark, major ($65.75\%$) and minor ($22.16\%$) triad qualities (Head~0, including extended chords) occupy $87.91\%$ of non-silent duration ($73.7\,\text{h}$, $11{,}436{,}357$ frames). Rare triads ({\small\texttt{sus4}}, {\small\texttt{dim}}, {\small\texttt{sus2}}, {\small\texttt{aug}}) collectively span under $3.5\%$ of duration, while altered upper extensions ({\small\texttt{\#11}} at $0.13\%$, {\small\texttt{b9}} at $0.13\%$, {\small\texttt{b13}} at $0.08\%$) inhabit a sparse tail. Moreover, human inter-annotator disagreement exceeds $40\%$ on seventh chords, tetrads, and their inversions~\cite{koops2019annotator}, compounding sample scarcity with label noise.

Recent architectures address large-vocabulary recognition via structured multi-task losses~\cite{mcfee2017structured}, 6-head chord decomposition~\cite{jiang2019large,akram2026chordformer}, and contrastive pseudo-labeling~\cite{li2025large}. Specifically, Jiang et al.~\cite{jiang2019large} and Akram et al.~\cite{akram2026chordformer} factorize chords into six categorical heads (root+triad, bass, and $7\text{th}, 9\text{th}, 11\text{th}, 13\text{th}$ degrees) trained under inverse-frequency reweighted cross-entropy. We hypothesize that loss reweighting alone cannot overcome a scarcity of physical acoustic exemplars: when an altered degree occurs in only a few dozen frames across a corpus, scaling its gradient penalty may amplify gradient variance and overfit to the instrumentation and room acoustics of those few tracks. Directly addressing the long tail motivates intervening on the training audio distribution.

We propose a Sim-to-Real pretraining framework driven by procedurally synthesized multitrack audio, contributing: (1) \textbf{Probabilistic Composition Generator}, an open-source engine sampling progressions along empirical transitions curated via the ChoCo aggregator~\cite{deberardinis2023choco} under deterministic quota scheduling (upweighting rare triads up to $15.6\times$ and altered extensions up to $12.3\times$), enforcing physical register bounds and $100\%$ extension sounding across four stems with random-mix augmentation to render $5{,}000$ tracks ($460.3\,\text{h}$) and frame-exact Harte labels~\cite{harte2005symbolic}; (2) \textbf{Two-Stage Sim-to-Real Curriculum}, a transfer regimen for a modified ChordFormer variant~\cite{akram2026chordformer} combining synthetic pretraining with real-audio finetuning via inverted discriminative learning rates ($\mathbf{s}_{\text{lr}} = [0.2, 0.5, 1.0]$) and floor curriculum replay ($\rho(\eta) = 0.150 \!\to\! 0.100$); and (3) \textbf{Multi-Head and WCSR Benchmark Analysis}, a 5-fold cross-validated evaluation over all $11{,}436{,}357$ frames showing recall improvements over the real-only baseline on targeted rare classes ($+6.64\,\text{pp}$ on rare triads, $+25.16\,\text{pp}$ on natural 9ths, $+5.46\,\text{pp}$ on {\small\texttt{\#9}}, $+5.74\,\text{pp}$ on {\small\texttt{b13}}) alongside an analysis of precision-recall trade-offs, duration-weighted WCSR bias, and polyphonic overtone masking.
